\documentclass[10pt]{article}
\usepackage{amsmath,amssymb}
\usepackage{longtable}

\usepackage[a4paper,left=2cm,right=2cm,top=1.5cm,bottom=2cm]{geometry}
\usepackage{xcolor}
\setlength{\parindent}{0pt}
\setlength{\emergencystretch}{2em}
\begin{document}

\title{Neural Networks by Hand}
\author{}
\date{}
\maketitle
\subsection*{Learning goals}
After this exercise, you should be able to

\begin{itemize}
\item
evaluate a simple neural network by hand,

\item
explain the role of weights, biases, and activation functions,

\item
perform a complete forward pass through a small neural network,

\item
calculate the mean squared error for a small training dataset,

\item
use backpropagation to compute gradients of the loss with respect to all trainable parameters,

\item
explain how the ReLU activation affects gradient propagation,

\item
perform one gradient-descent update,

\item
verify that a gradient step can reduce the training loss,

\item
distinguish between fitting the training points and learning the underlying function,

\item
explain why a ReLU network produces a piecewise-linear function.

\end{itemize}
The calculations in this sheet are designed to be done by hand.

\clearpage
\section*{Part I - A small neural network}
We want to approximate the function

\[
f(x)=-x^2+4x
\]

We use only the four training inputs

\[
x_i=0,\ 1,\ 2,\ 3.
\]

\subsection*{Exercise 1.1 - Generate the training data}
Evaluate the function at the four training inputs and complete the table.

\begingroup
\small
\setlength{\tabcolsep}{4pt}
\begin{longtable}{l|l}
$x_i$ & $y_i=f(x_i)$ \\[1.5pt]
\hline
\endfirsthead
$x_i$ & $y_i=f(x_i)$ \\[1.5pt]
\hline
\endhead
\rule{0pt}{2.4ex}0 &  $\phantom{-0.00}$ \\[3pt]
\rule{0pt}{2.4ex}1 &  $\phantom{-0.00}$ \\[3pt]
\rule{0pt}{2.4ex}2 &  $\phantom{-0.00}$ \\[3pt]
\rule{0pt}{2.4ex}3 &  $\phantom{-0.00}$ \\[3pt]
\end{longtable}
\endgroup
Plot the four training points and sketch the true function $f(x)=-x^2+4x$ in the range $-1\leq x\leq5$.

\subsection*{Exercise 1.2 - Our neural network architecture}
We use the architecture

\[
1\rightarrow2\rightarrow1.
\]

There is

\begin{itemize}
\item
one input

\item
one hidden layer with two neurons

\item
one output neuron

\end{itemize}
The hidden layer uses the $ReLU$ activation

\[
\operatorname{ReLU}(z)=\max(0,z)
\]

The two hidden neurons are

\[
z_{1,i}=w_1x_i+b_1
\]

\[
h_{1,i}=\operatorname{ReLU}(z_{1,i})
\]

and

\[
z_{2,i}=w_2x_i+b_2
\]

\[
h_{2,i}=\operatorname{ReLU}(z_{2,i})
\]

The output layer is linear:

\[
\hat y_i=v_1h_{1,i}+v_2h_{2,i}+c
\]

Sketch the following

\begin{itemize}
\item
the $ReLU$ function for input values $z\in[-5,5]$

\item
the neural network

\end{itemize}
for this, count the trainable parameters first. Your drawing should contain

\begin{itemize}
\item
the input $x_i$

\item
the two hidden neurons

\item
the output $\hat y_i$

\item
all weights

\item
all biases

\item
the ReLU activation in the hidden layer

\end{itemize}
\subsection*{Exercise 1.3 - A complete forward pass}
To build the neural network, the parameters need to be fit. We start with a list of initial guesses for which the output of the neural network is calculated. This will be compared to the training data using the mean squared error. A real neural network training proceeds until the MSE is sufficiently small. In this exercise we will however only produce a single calculation.

For the first training step, use the following initial parameters:

\[
w_1=1
\qquad
b_1=-0.5
\]

\[
w_2=-1
\qquad
b_2=2.5
\]

\[
v_1=1
\qquad
v_2=1
\qquad
c=0
\]

This exercise is divised into multiple steps:

\subsubsection*{Step 1 - Hidden-layer activations}
For each of the four training inputs, calculate

\[
z_{1,i}=w_1x_i+b_1
\]

\[
h_{1,i}=\operatorname{ReLU}(z_{1,i})
\]

and analoguous for $z_{2,i}$ and $h_{2,i}$. Complete the table below:

\begingroup
\small
\setlength{\tabcolsep}{4pt}
\begin{longtable}{l|l|l|l|l|l}
$x_i$ & $y_i$ & $z_{1,i}$ & $h_{1,i}$ & $z_{2,i}$ & $h_{2,i}$ \\[1.5pt]
\hline
\endfirsthead
$x_i$ & $y_i$ & $z_{1,i}$ & $h_{1,i}$ & $z_{2,i}$ & $h_{2,i}$ \\[1.5pt]
\hline
\endhead
\rule{0pt}{2.4ex}0 & 0 &  $\phantom{-0.00}$ &  $\phantom{-0.00}$ &  $\phantom{-0.00}$ &  $\phantom{-0.00}$ \\[3pt]
\rule{0pt}{2.4ex}1 & 3 &  $\phantom{-0.00}$ &  $\phantom{-0.00}$ &  $\phantom{-0.00}$ &  $\phantom{-0.00}$ \\[3pt]
\rule{0pt}{2.4ex}2 & 4 &  $\phantom{-0.00}$ &  $\phantom{-0.00}$ &  $\phantom{-0.00}$ &  $\phantom{-0.00}$ \\[3pt]
\rule{0pt}{2.4ex}3 & 3 &  $\phantom{-0.00}$ &  $\phantom{-0.00}$ &  $\phantom{-0.00}$ &  $\phantom{-0.00}$ \\[3pt]
\end{longtable}
\endgroup
For each hidden neuron, identify for which input values the ReLU is active.

\subsubsection*{Step 2 - Output predictions}
Using

\[
\hat y=v_1h_1+v_2h_2+c
\]

calculate the network prediction for every training point:


\begin{center}
\begingroup
\small
\setlength{\tabcolsep}{4pt}
\begin{tabular}{l|l|l|l}
$x$ & $y$ & $\hat y$ & Residual $\hat y-y$ \\ \hline
\rule{0pt}{2.4ex}0 & 0 & $\phantom{0}$ & $\phantom{0}$ \\[3pt]
\rule{0pt}{2.4ex}1 & 3 & $\phantom{0}$ & $\phantom{0}$ \\[3pt]
\rule{0pt}{2.4ex}2 & 4 & $\phantom{0}$ & $\phantom{0}$ \\[3pt]
\rule{0pt}{2.4ex}3 & 3 & $\phantom{0}$ & $\phantom{0}$ \\[3pt]
\end{tabular}
\endgroup
\end{center}


Plot the four predictions together with the four training points.

\subsubsection*{Step 3 - Calculate the loss}
We use the mean squared error

\[
L=\frac{1}{N}\sum_{i=1}^{N}(\hat y_i-y_i)^2
\]

with

\[
N=4
\]

Calculate the MSE for the current network. This is the loss before training.

\subsection*{Exercise 1.4 - Backpropagation by hand}
We now calculate how the loss changes with respect to every trainable parameter.

The important idea is:

\[
\text{"Start at the loss and repeatedly apply the chain rule while moving backwards through the network"}
\]

We perform this calculation only once by hand. A neural-network library repeats exactly the same type of calculation automatically.

\subsubsection*{Step 1 - Start at the loss}
For one training sample the loss is

\[
L_i=(\hat y_i-y_i)^2
\]

and so the mean loss is given by a sum

\[
L=\frac{1}{4}\sum_i L_i
\]

The gradient is given by

\[
\frac{\partial L}{\partial \hat y_i}=\frac{2}{4}(\hat y_i-y_i)
\]

One defines

\[
\delta_i=\frac{\partial L}{\partial \hat y_i}
\]

Calculate $\delta_i$ for every training point.

\begingroup
\small
\setlength{\tabcolsep}{4pt}
\begin{longtable}{l|l|l|l|l}
$x_i$ & $y_i$ & $\hat y_i$ & $\hat y_i-y_i$ & $\delta=\partial L/\partial\hat y$ \\[1.5pt]
\hline
\endfirsthead
$x_i$ & $y_i$ & $\hat y_i$ & $\hat y_i-y_i$ & $\delta=\partial L/\partial\hat y$ \\[1.5pt]
\hline
\endhead
\rule{0pt}{2.4ex}0 & 0 &  $\phantom{-0.00}$ &  $\phantom{-0.00}$ &  $\phantom{-0.00}$ \\[3pt]
\rule{0pt}{2.4ex}1 & 3 &  $\phantom{-0.00}$ &  $\phantom{-0.00}$ &  $\phantom{-0.00}$ \\[3pt]
\rule{0pt}{2.4ex}2 & 4 &  $\phantom{-0.00}$ &  $\phantom{-0.00}$ &  $\phantom{-0.00}$ \\[3pt]
\rule{0pt}{2.4ex}3 & 3 &  $\phantom{-0.00}$ &  $\phantom{-0.00}$ &  $\phantom{-0.00}$ \\[3pt]
\end{longtable}
\endgroup
\subsubsection*{Step 2 - Gradients of the output layer}
\paragraph*{Gradient with respect to $v_1$}
The output is

\[
\hat y_i=v_1h_{1,i}+v_2h_{2,i}+c
\]

thus

\[
\frac{\partial \hat y_i}{\partial v_1}=h_{1,i}
\]

Use the chain rule to obtain

\[
\frac{\partial L}{\partial v_1}=\sum_i\frac{\partial L}{\partial\hat y_i}\frac{\partial\hat y_i}{\partial v_1}
\]

Therefore,

\[
\frac{\partial L}{\partial v_1}=\sum_i \delta_i h_{1,i}
\]

Calculate this gradient.

\paragraph*{Gradient with respect to $v_2$}
Analogously, calculate

\[
\frac{\partial L}{\partial v_2}=\sum_i \delta_i h_{2,i}
\]

\paragraph*{Gradient with respect to $c$}
Since

\[
\frac{\partial\hat y_i}{\partial c}=1
\]

calculate

\[
\frac{\partial L}{\partial c}=\sum_i\delta_i
\]

Record your results:

\begingroup
\small
\setlength{\tabcolsep}{4pt}
\begin{longtable}{l|l}
parameter & gradient \\[1.5pt]
\hline
\endfirsthead
parameter & gradient \\[1.5pt]
\hline
\endhead
\rule{0pt}{2.4ex}$v_1$ &  $\phantom{-0.00}$ \\[3pt]
\rule{0pt}{2.4ex}$v_2$ &  $\phantom{-0.00}$ \\[3pt]
\rule{0pt}{2.4ex}$c$ &  $\phantom{-0.00}$ \\[3pt]
\end{longtable}
\endgroup
\subsubsection*{Step 3 - Backpropagate through ReLU}
The derivative of ReLU is

\[
\operatorname{ReLU}'(z)=
\begin{cases}
0, & z<0,\\1, & z>0
\end{cases}
\]

For the values encountered in this exercise, none of the $z$ values are exactly zero.

Using your values from the forward pass, complete the table:

\begingroup
\small
\setlength{\tabcolsep}{4pt}
\begin{longtable}{l|l|l|l|l}
$x_i$ & $z_{1,i}$ & $\operatorname{ReLU}'(z_{1,i})$ & $z_{2,i}$ & $\operatorname{ReLU}'(z_{2,i})$ \\[1.5pt]
\hline
\endfirsthead
$x_i$ & $z_{1,i}$ & $\operatorname{ReLU}'(z_{1,i})$ & $z_{2,i}$ & $\operatorname{ReLU}'(z_{2,i})$ \\[1.5pt]
\hline
\endhead
\rule{0pt}{2.4ex}0 &  $\phantom{-0.00}$ &  $\phantom{-0.00}$ &  $\phantom{-0.00}$ &  $\phantom{-0.00}$ \\[3pt]
\rule{0pt}{2.4ex}1 &  $\phantom{-0.00}$ &  $\phantom{-0.00}$ &  $\phantom{-0.00}$ &  $\phantom{-0.00}$ \\[3pt]
\rule{0pt}{2.4ex}2 &  $\phantom{-0.00}$ &  $\phantom{-0.00}$ &  $\phantom{-0.00}$ &  $\phantom{-0.00}$ \\[3pt]
\rule{0pt}{2.4ex}3 &  $\phantom{-0.00}$ &  $\phantom{-0.00}$ &  $\phantom{-0.00}$ &  $\phantom{-0.00}$ \\[3pt]
\end{longtable}
\endgroup
For the first hidden neuron,

\[
\frac{\partial L}{\partial z_{1,i}}=\frac{\partial L}{\partial\hat y_i}\frac{\partial\hat y_i}{\partial h_{1,i}}\frac{\partial h_{1,i}}{\partial z_{1,i}}
\]

Because

\[
\frac{\partial\hat y_i}{\partial h_{1,i}}=v_1
\]

we obtain

\[
\frac{\partial L}{\partial z_{1,i}}=\delta_i v_1\operatorname{ReLU}'(z_{1,i})
\]

Calculate these four values.

Do the same for the second hidden neuron:

\[
\frac{\partial L}{\partial z_{2,i}}=\delta_i v_2\operatorname{ReLU}'(z_{2,i})
\]

Complete:

\begingroup
\small
\setlength{\tabcolsep}{4pt}
\begin{longtable}{l|l|l}
$x_i$ & $\partial L/\partial z_{1,i}$ & $\partial L/\partial z_{2,i}$ \\[1.5pt]
\hline
\endfirsthead
$x_i$ & $\partial L/\partial z_{1,i}$ & $\partial L/\partial z_{2,i}$ \\[1.5pt]
\hline
\endhead
\rule{0pt}{2.4ex}0 &  $\phantom{-0.00}$ &  $\phantom{-0.00}$ \\[3pt]
\rule{0pt}{2.4ex}1 &  $\phantom{-0.00}$ &  $\phantom{-0.00}$ \\[3pt]
\rule{0pt}{2.4ex}2 &  $\phantom{-0.00}$ &  $\phantom{-0.00}$ \\[3pt]
\rule{0pt}{2.4ex}3 &  $\phantom{-0.00}$ &  $\phantom{-0.00}$ \\[3pt]
\end{longtable}
\endgroup
What happens to the gradient when a ReLU neuron is inactive?

\subsubsection*{Step 4 - Gradients of the first layer}
For the first neuron,

\[
z_{1,i}=w_1x_i+b_1.
\]

Therefore,

\[
\frac{\partial z_{1,i}}{\partial w_1}=x_i
\]

and

\[
\frac{\partial z_{1,i}}{\partial b_1}=1
\]

Calculate the gradients

\[
\frac{\partial L}{\partial w_1}
\]

and

\[
\frac{\partial L}{\partial b_1}
\]

Repeat the calculation for the second hidden neuron and collect all seven gradients:

\begingroup
\small
\setlength{\tabcolsep}{4pt}
\begin{longtable}{l|l}
parameter & gradient \\[1.5pt]
\hline
\endfirsthead
parameter & gradient \\[1.5pt]
\hline
\endhead
\rule{0pt}{2.4ex}$w_1$ &  $\phantom{-0.00}$ \\[3pt]
\rule{0pt}{2.4ex}$b_1$ &  $\phantom{-0.00}$ \\[3pt]
\rule{0pt}{2.4ex}$w_2$ &  $\phantom{-0.00}$ \\[3pt]
\rule{0pt}{2.4ex}$b_2$ &  $\phantom{-0.00}$ \\[3pt]
\rule{0pt}{2.4ex}$v_1$ &  $\phantom{-0.00}$ \\[3pt]
\rule{0pt}{2.4ex}$v_2$ &  $\phantom{-0.00}$ \\[3pt]
\rule{0pt}{2.4ex}$c$ &  $\phantom{-0.00}$ \\[3pt]
\end{longtable}
\endgroup
\subsection*{Exercise 1.5 - One gradient-descent step}
We now update every parameter using gradient descent. For a parameter $\theta$

\[
\theta_{\mathrm{new}}=\theta_{\mathrm{old}}-\eta\frac{\partial L}{\partial\theta}
\]

Use the learning rate

\[
\eta=0.1
\]

\subsubsection*{Step 1 - Update the parameters}
Complete the table.

\begingroup
\small
\setlength{\tabcolsep}{4pt}
\begin{longtable}{l|l|l|l}
parameter & old value & gradient & new value \\[1.5pt]
\hline
\endfirsthead
parameter & old value & gradient & new value \\[1.5pt]
\hline
\endhead
\rule{0pt}{2.4ex}$w_1$ & 1.0 &  $\phantom{-0.00}$ &  $\phantom{-0.00}$ \\[3pt]
\rule{0pt}{2.4ex}$b_1$ & -0.5 &  $\phantom{-0.00}$ &  $\phantom{-0.00}$ \\[3pt]
\rule{0pt}{2.4ex}$w_2$ & -1.0 &  $\phantom{-0.00}$ &  $\phantom{-0.00}$ \\[3pt]
\rule{0pt}{2.4ex}$b_2$ & 2.5 &  $\phantom{-0.00}$ &  $\phantom{-0.00}$ \\[3pt]
\rule{0pt}{2.4ex}$v_1$ & 1.0 &  $\phantom{-0.00}$ &  $\phantom{-0.00}$ \\[3pt]
\rule{0pt}{2.4ex}$v_2$ & 1.0 &  $\phantom{-0.00}$ &  $\phantom{-0.00}$ \\[3pt]
\rule{0pt}{2.4ex}$c$ & 0.0 &  $\phantom{-0.00}$ &  $\phantom{-0.00}$ \\[3pt]
\end{longtable}
\endgroup
\subsubsection*{Step 2 - Did training improve the model?}
Using the updated parameters, perform one more forward pass (see Exercise 1.3). (You do not need to perform another backward pass)

Calculate the new predictions and the new MSE.

\begingroup
\small
\setlength{\tabcolsep}{4pt}
\begin{longtable}{l|l|l}
$x$ & $y$ & new prediction $\hat y_{\mathrm{new}}$ \\[1.5pt]
\hline
\endfirsthead
$x$ & $y$ & new prediction $\hat y_{\mathrm{new}}$ \\[1.5pt]
\hline
\endhead
\rule{0pt}{2.4ex}0 & 0 &  $\phantom{-0.00}$ \\[3pt]
\rule{0pt}{2.4ex}1 & 3 &  $\phantom{-0.00}$ \\[3pt]
\rule{0pt}{2.4ex}2 & 4 &  $\phantom{-0.00}$ \\[3pt]
\rule{0pt}{2.4ex}3 & 3 &  $\phantom{-0.00}$ \\[3pt]
\end{longtable}
\endgroup
Compare

\[
L_{\mathrm{old}}
\]

and

\[
L_{\mathrm{new}}.
\]

Did this gradient-descent step reduce the loss? Does every single training point necessarily improve after one gradient step?

Explain why the optimizer can still reduce the total loss even if one individual prediction becomes worse.

\clearpage
\section*{Part II - Investigating a fitted neural network}
Repeating forward pass, backpropagation, and parameter updates many times is exactly what neural-network training does.

Instead of repeating these calculations by hand, we now inspect a network whose parameters are already chosen.

Consider

\[
h_1(x)=\operatorname{ReLU}(x),
\]

\[
h_2(x)=\operatorname{ReLU}(x-1.5),
\]

and

\[
\hat f(x)=3h_1(x)-4h_2(x).
\]

This is again a network with architecture

\[
1\rightarrow2\rightarrow1.
\]

\subsection*{Exercise 2.1 - Identify the network parameters}
Write the network in the standard form

\[
z_1=w_1x+b_1,
\qquad
h_1=\operatorname{ReLU}(z_1),
\]

\[
z_2=w_2x+b_2,
\qquad
h_2=\operatorname{ReLU}(z_2),
\]

\[
\hat y=v_1h_1+v_2h_2+c.
\]

Determine the parameters

\[
w_1,\ b_1,\ w_2,\ b_2,\ v_1,\ v_2,\ c.
\]

\subsection*{Exercise 2.2 - Check the training points}
Evaluate the fitted network at

\[
x=0,\ 1,\ 2,\ 3.
\]

Complete the table.

\begingroup
\small
\setlength{\tabcolsep}{4pt}
\begin{longtable}{l|l|l}
$x$ & true $f(x)$ & network $\hat f(x)$ \\[1.5pt]
\hline
\endfirsthead
$x$ & true $f(x)$ & network $\hat f(x)$ \\[1.5pt]
\hline
\endhead
\rule{0pt}{2.4ex}0 & 0 &  $\phantom{-0.00}$ \\[3pt]
\rule{0pt}{2.4ex}1 & 3 &  $\phantom{-0.00}$ \\[3pt]
\rule{0pt}{2.4ex}2 & 4 &  $\phantom{-0.00}$ \\[3pt]
\rule{0pt}{2.4ex}3 & 3 &  $\phantom{-0.00}$ \\[3pt]
\end{longtable}
\endgroup
Calculate the training MSE.

What is special about this result?

\subsection*{Exercise 2.3 - Has the network learned the true function?}
Now evaluate both the true function and the neural network at new input values:

\[
x=0.5,\ 1.5,\ 2.5,\ 4.
\]

Complete the table.

\begingroup
\small
\setlength{\tabcolsep}{4pt}
\begin{longtable}{l|l|l|l}
$x$ & true $f(x)$ & network $\hat f(x)$ & difference \\[1.5pt]
\hline
\endfirsthead
$x$ & true $f(x)$ & network $\hat f(x)$ & difference \\[1.5pt]
\hline
\endhead
\rule{0pt}{2.4ex}0.5 &  $\phantom{-0.00}$ &  $\phantom{-0.00}$ &  $\phantom{-0.00}$ \\[3pt]
\rule{0pt}{2.4ex}1.5 &  $\phantom{-0.00}$ &  $\phantom{-0.00}$ &  $\phantom{-0.00}$ \\[3pt]
\rule{0pt}{2.4ex}2.5 &  $\phantom{-0.00}$ &  $\phantom{-0.00}$ &  $\phantom{-0.00}$ \\[3pt]
\rule{0pt}{2.4ex}4.0 &  $\phantom{-0.00}$ &  $\phantom{-0.00}$ &  $\phantom{-0.00}$ \\[3pt]
\end{longtable}
\endgroup
Discuss:

\begin{enumerate}
\item
The network has zero training error. Has it learned the exact function?

\item
Why is performance on unseen inputs important?

\end{enumerate}
\subsection*{Exercise 2.4 - What function has the ReLU network learned?}
The fitted network is

\[
\hat f(x)=3\operatorname{ReLU}(x)-4\operatorname{ReLU}(x-1.5)
\]

Because ReLU changes its behavior at zero, the expression changes at

\[
x=0
\]

and

\[
x=1.5
\]

Consider the three regions separately and determine $\hat f(x)$. Write the final result as

\[
\hat f(x)=
\begin{cases}
\ldots, & x<0,\\
\ldots, & 0\leq x<1.5,\\
\ldots, & x\geq1.5
\end{cases}
\]

Plot the function. What is the most important visual difference between the quadratic target function and the ReLU network?

\clearpage
\section*{Part III - A deeper neural network}
So far, we have considered a neural network with one hidden layer. We now investigate a deeper network with the architecture

\[
1\rightarrow3\rightarrow3\rightarrow1
\]

There is

\begin{itemize}
\item
one input

\item
a first hidden layer with three neurons

\item
a second hidden layer with three neurons

\item
one output neuron

\end{itemize}
Both hidden layers use the $ReLU$ activation function.

In this part, the parameters are already given. You do not need to perform backpropagation for this larger network. The goal is to understand how information is propagated through multiple hidden layers and how additional neurons can change the shape of the learned function.

\subsection*{Exercise 3.1 - Architecture and number of parameters}
For a fully connected neural network, every neuron in one layer is connected to every neuron in the next layer.

\begin{enumerate}
\item
Draw the architecture

\end{enumerate}
\[
1\rightarrow3\rightarrow3\rightarrow1
\]

\begin{enumerate}
\setcounter{enumi}{1}
\item
Count the number of weights and biases between the input and the first hidden layer.

\item
Count the number of weights and biases between the first and second hidden layers.

\item
Count the number of weights and biases between the second hidden layer and the output.

\item
Determine the total number of trainable parameters.

\end{enumerate}
Compare this number with the seven trainable parameters of the network from Parts I and II.

\subsection*{Exercise 3.2 - First hidden layer}
The first hidden layer is defined by

\[
z^{(1)}_{1,i}=x_i,
\qquad
h^{(1)}_{1,i}=\operatorname{ReLU}(z^{(1)}_{1,i}),
\]

\[
z^{(1)}_{2,i}=x_i-1,
\qquad
h^{(1)}_{2,i}=\operatorname{ReLU}(z^{(1)}_{2,i}),
\]

and

\[
z^{(1)}_{3,i}=x_i-2,
\qquad
h^{(1)}_{3,i}=\operatorname{ReLU}(z^{(1)}_{3,i}).
\]

Calculate the activations for the four training inputs.

\begingroup
\small
\setlength{\tabcolsep}{4pt}
\begin{longtable}{l|l|l|l|l|l|l}
$x_i$ & $z^{(1)}_{1,i}$ & $h^{(1)}_{1,i}$ & $z^{(1)}_{2,i}$ & $h^{(1)}_{2,i}$ & $z^{(1)}_{3,i}$ & $h^{(1)}_{3,i}$ \\[1.5pt]
\hline
\endfirsthead
$x_i$ & $z^{(1)}_{1,i}$ & $h^{(1)}_{1,i}$ & $z^{(1)}_{2,i}$ & $h^{(1)}_{2,i}$ & $z^{(1)}_{3,i}$ & $h^{(1)}_{3,i}$ \\[1.5pt]
\hline
\endhead
\rule{0pt}{2.4ex}0 &  $\phantom{-0.00}$ &  $\phantom{-0.00}$ &  $\phantom{-0.00}$ &  $\phantom{-0.00}$ &  $\phantom{-0.00}$ &  $\phantom{-0.00}$ \\[3pt]
\rule{0pt}{2.4ex}1 &  $\phantom{-0.00}$ &  $\phantom{-0.00}$ &  $\phantom{-0.00}$ &  $\phantom{-0.00}$ &  $\phantom{-0.00}$ &  $\phantom{-0.00}$ \\[3pt]
\rule{0pt}{2.4ex}2 &  $\phantom{-0.00}$ &  $\phantom{-0.00}$ &  $\phantom{-0.00}$ &  $\phantom{-0.00}$ &  $\phantom{-0.00}$ &  $\phantom{-0.00}$ \\[3pt]
\rule{0pt}{2.4ex}3 &  $\phantom{-0.00}$ &  $\phantom{-0.00}$ &  $\phantom{-0.00}$ &  $\phantom{-0.00}$ &  $\phantom{-0.00}$ &  $\phantom{-0.00}$ \\[3pt]
\end{longtable}
\endgroup
At which input value does each of the three neurons become active?

How does this differ from the two-neuron hidden layer considered previously?

\subsection*{Exercise 3.3 - Second hidden layer}
The second hidden layer receives the three outputs of the first hidden layer as its inputs.

It is defined by

\[
z^{(2)}_{1,i}=h^{(1)}_{1,i},
\qquad
h^{(2)}_{1,i}=\operatorname{ReLU}(z^{(2)}_{1,i})
\]

\[
z^{(2)}_{2,i}=h^{(1)}_{1,i}-h^{(1)}_{2,i},
\qquad
h^{(2)}_{2,i}=\operatorname{ReLU}(z^{(2)}_{2,i})
\]

and

\[
z^{(2)}_{3,i}=h^{(1)}_{2,i}-h^{(1)}_{3,i},
\qquad
h^{(2)}_{3,i}=\operatorname{ReLU}(z^{(2)}_{3,i})
\]

Use the results from Exercise 3.2 and complete the table.

\begingroup
\small
\setlength{\tabcolsep}{4pt}
\begin{longtable}{l|l|l|l|l|l|l}
$x$ & $z^{(2)}_{1,i}$ & $h^{(2)}_{1,i}$ & $z^{(2)}_{2,i}$ & $h^{(2)}_{2,i}$ & $z^{(2)}_{3,i}$ & $h^{(2)}_{3,i}$ \\[1.5pt]
\hline
\endfirsthead
$x$ & $z^{(2)}_{1,i}$ & $h^{(2)}_{1,i}$ & $z^{(2)}_{2,i}$ & $h^{(2)}_{2,i}$ & $z^{(2)}_{3,i}$ & $h^{(2)}_{3,i}$ \\[1.5pt]
\hline
\endhead
\rule{0pt}{2.4ex}0 &  $\phantom{-0.00}$ &  $\phantom{-0.00}$ &  $\phantom{-0.00}$ &  $\phantom{-0.00}$ &  $\phantom{-0.00}$ &  $\phantom{-0.00}$ \\[3pt]
\rule{0pt}{2.4ex}1 &  $\phantom{-0.00}$ &  $\phantom{-0.00}$ &  $\phantom{-0.00}$ &  $\phantom{-0.00}$ &  $\phantom{-0.00}$ &  $\phantom{-0.00}$ \\[3pt]
\rule{0pt}{2.4ex}2 &  $\phantom{-0.00}$ &  $\phantom{-0.00}$ &  $\phantom{-0.00}$ &  $\phantom{-0.00}$ &  $\phantom{-0.00}$ &  $\phantom{-0.00}$ \\[3pt]
\rule{0pt}{2.4ex}3 &  $\phantom{-0.00}$ &  $\phantom{-0.00}$ &  $\phantom{-0.00}$ &  $\phantom{-0.00}$ &  $\phantom{-0.00}$ &  $\phantom{-0.00}$ \\[3pt]
\end{longtable}
\endgroup
Notice that the second hidden layer no longer acts directly on $x$. It transforms features that were already created by the first hidden layer.

\subsection*{Exercise 3.4 - Output of the deep network}
The output layer is

\[
\hat f_{\mathrm{deep}}(x)=-h^{(2)}_1+4h^{(2)}_2+2h^{(2)}_3.
\]

Calculate the prediction for every training point.

\begingroup
\small
\setlength{\tabcolsep}{4pt}
\begin{longtable}{l|l|l}
$x_i$ & true $f(x_i)$ & deep network $\hat f_{\mathrm{deep}}(x_i)$ \\[1.5pt]
\hline
\endfirsthead
$x_i$ & true $f(x_i)$ & deep network $\hat f_{\mathrm{deep}}(x_i)$ \\[1.5pt]
\hline
\endhead
\rule{0pt}{2.4ex}0 & 0 &  $\phantom{-0.00}$ \\[3pt]
\rule{0pt}{2.4ex}1 & 3 &  $\phantom{-0.00}$ \\[3pt]
\rule{0pt}{2.4ex}2 & 4 &  $\phantom{-0.00}$ \\[3pt]
\rule{0pt}{2.4ex}3 & 3 &  $\phantom{-0.00}$ \\[3pt]
\end{longtable}
\endgroup
Calculate the training MSE.

Compare the result with the fitted $1\rightarrow2\rightarrow1$ network from Part II.

\begin{itemize}
\item
Do both networks achieve the same training error?

\item
Does this imply that both networks represent the same function?

\end{itemize}
\subsection*{Exercise 3.5 - Test the network between the training points}
The two networks both reproduce the four training points exactly. We now compare what they predict between the training points.

Evaluate

\[
f(x)=-x^2+4x
\]

at

\[
x=0.5,\ 1.25,\ 1.75,\ 2.5
\]

and use your result from the shallow network of Part II. For the deep network, perform the forward pass through both hidden layers and complete the table.

\begingroup
\small
\setlength{\tabcolsep}{4pt}
\begin{longtable}{l|l|l|l}
$x_i$ & true $f(x_i)$ & shallow network & deep network \\[1.5pt]
\hline
\endfirsthead
$x_i$ & true $f(x_i)$ & shallow network & deep network \\[1.5pt]
\hline
\endhead
\rule{0pt}{2.4ex}0.5 &  $\phantom{-0.00}$ &  $\phantom{-0.00}$ &  $\phantom{-0.00}$ \\[3pt]
\rule{0pt}{2.4ex}1.25 &  $\phantom{-0.00}$ &  $\phantom{-0.00}$ &  $\phantom{-0.00}$ \\[3pt]
\rule{0pt}{2.4ex}1.75 &  $\phantom{-0.00}$ &  $\phantom{-0.00}$ &  $\phantom{-0.00}$ \\[3pt]
\rule{0pt}{2.4ex}2.5 &  $\phantom{-0.00}$ &  $\phantom{-0.00}$ &  $\phantom{-0.00}$ \\[3pt]
\end{longtable}
\endgroup
Which network is closer to the true function at these points? Why can two networks with zero training error behave differently between the training observations?

\subsection*{Exercise 3.6 - What function has the deeper network learned?}
Determine the output of the deeper network in the following regions:

\begin{enumerate}
\item
$x<0$

\item
$0\leq x<1$

\item
$1\leq x<2$

\item
$x\geq2$

\end{enumerate}
Write the result as a piecewise function

\[
\hat f_{\mathrm{deep}}(x)=
\begin{cases}
\ldots, & x<0,\\
\ldots, & 0\leq x<1,\\
\ldots, & 1\leq x<2,\\
\ldots, & x\geq2
\end{cases}
\]

Plot in the same coordinate system

\begin{itemize}
\item
the true function $f(x)=-x^2+4x$

\item
the four training points

\item
the fitted $1\rightarrow2\rightarrow1$ network from Part II

\item
the fitted $1\rightarrow3\rightarrow3\rightarrow1$ network

\end{itemize}
What has changed when moving to the larger network?

In particular, compare the number and positions of the changes in slope.

\clearpage
\section*{Part IV - Short conceptual questions}
\subsection*{Question 1}
Why is the output of this ReLU network piecewise linear?

\subsection*{Question 2}
What determines the positions at which the slope changes?

\subsection*{Question 3}
What role do the output weights $v_1$ and $v_2$ play?

\subsection*{Question 4}
Why can a network with only two hidden neurons already represent a function that is more complicated than a single straight line?

\subsection*{Question 5}
What would additional ReLU neurons allow the network to do?

\subsection*{Question 6}
Why does zero training error not guarantee good generalization?


\end{document}
